% Appendix
% BSc Thesis: Improving RAG Systems for Earth Observation
% Through Multi-Agent Ontology-Based Retrieval
% Author: Pieterjan Pauwelyn

\chapter{Problem Analysis Questions}\label{app:analysis-questions}

This appendix lists the 30 domain-specific questions used for the manual
problem analysis in Chapter~\ref{chap:preliminaries}. Twenty questions
span five intent categories; the remaining ten are highly technical
questions that do not fit a single category.

\medskip
\noindent\textbf{Exploratory questions}
\begin{itemize}
\item ``What satellite sensors can be used to monitor vegetation health?''
\item ``How can synthetic aperture radar detect selective logging activities in tropical forests during the rainy season?''
\item ``What combination of lidar metrics and hyperspectral indices best characterizes understory vegetation diversity in temperate mixed forests?''
\end{itemize}

\noindent\textbf{Comparative questions}
\begin{itemize}
\item ``How do optical and radar sensors compare for forest monitoring?''
\item ``How does the accuracy of above-ground biomass estimation compare between GEDI lidar data and Sentinel-1 SAR backscatter in boreal forests?''
\item ``Compare the performance of random forest versus deep learning approaches for tree species classification using WorldView-3 imagery in mixed deciduous forests.''
\end{itemize}

\noindent\textbf{Descriptive questions}
\begin{itemize}
\item ``What is NDVI and how is it used in vegetation monitoring?''
\item ``Explain the role of the Hansen Global Forest Change dataset in monitoring tropical deforestation patterns.''
\item ``Describe how full-waveform lidar data characterizes vertical forest structure for biodiversity assessment in old-growth forests.''
\end{itemize}

\noindent\textbf{Causal questions}
\begin{itemize}
\item ``Why does forest fire detection use thermal infrared bands?''
\item ``How do edge effects influence carbon stock estimates in fragmented forests when using medium-resolution satellite data?''
\item ``Why do L-band SAR systems show better correlation with forest biomass than C-band in high-biomass tropical forests, and what are the physical mechanisms behind signal saturation?''
\end{itemize}

\noindent\textbf{Relational questions}
\begin{itemize}
\item ``What is the relationship between canopy cover and land surface temperature?''
\item ``How does the timing of Sentinel-2 image acquisition relate to phenological stages for deciduous forest species mapping?''
\item ``What is the relationship between interferometric coherence from Sentinel-1 and forest disturbance severity in managed forests?''
\end{itemize}

\noindent\textbf{Additional domain-specific questions}
\begin{itemize}
\item ``What are the specific advantages of using Sentinel-2's red-edge bands for forest health assessment?''
\item ``How would you design a multi-sensor approach to map forest degradation in a tropical rainforest with persistent cloud cover?''
\item ``What are the limitations of using the Hansen Global Forest Change dataset for monitoring forest degradation in dry tropical forests?''
\item ``How does canopy height from GEDI lidar data relate to carbon sequestration estimates, and what are the uncertainties?''
\item ``What specific Earth observation data and processing steps would you need to create a monthly forest fire risk map for Mediterranean forests?''
\item ``How do phenological variations affect the accuracy of tree species classification using Sentinel-2, and how would you account for this?''
\item ``What is the minimum detectable change in forest canopy cover using Landsat time series analysis?''
\item ``How can Planet's daily imagery be integrated with Sentinel-2 for near real-time deforestation alerts?''
\item ``What are the key spectral and structural indicators for detecting forest insect outbreak damage using multispectral and SAR data?''
\item ``Explain how GLAS/ICESat-2 altimetry data can be used to validate canopy height models derived from optical stereo imagery.''
\item ``What role does texture analysis play in distinguishing primary forest from secondary regrowth in high-resolution satellite imagery?''
\item ``How can interferometric SAR coherence time series detect gradual forest degradation that optical sensors might miss?''
\item ``What are the optimal seasonal windows for acquiring satellite imagery to map invasive species in temperate forests?''
\item ``How do you account for topographic effects when calculating vegetation indices in mountainous forest areas?''
\item ``What is the relationship between SAR backscatter seasonal variation and forest water stress in Mediterranean ecosystems?''
\end{itemize}

\chapter{Evaluation Criteria Definitions}\label{app:eval-criteria}

Table~\ref{tab:eval-criteria-full} defines the eight criteria used by the
final LLM-as-a-judge evaluation framework (Claude Sonnet~4.6). Each criterion
is described by a guiding question and four score bands spanning the 1--10
scale. See Chapter~\ref{chap:evaluation} for the rationale behind each
criterion and its connection to the failure patterns identified in
Chapter~\ref{chap:preliminaries}.

\begin{longtable}{@{}p{2.6cm} p{4.0cm} p{1.2cm} p{6.0cm}@{}}
\caption{Evaluation criteria, guiding questions, and score-band
descriptions used in the final LLM-as-a-judge framework.}
\label{tab:eval-criteria-full} \\
\toprule
\textbf{Criterion} & \textbf{Guiding question} & \textbf{Band} & \textbf{Description} \\
\midrule
\endfirsthead

\multicolumn{4}{c}{\small\emph{Table~\ref{tab:eval-criteria-full} continued}} \\
\toprule
\textbf{Criterion} & \textbf{Guiding question} & \textbf{Band} & \textbf{Description} \\
\midrule
\endhead

\bottomrule
\endfoot

% scientific accuracy
\multirow{4}{2.6cm}{Scientific Accuracy}
& \multirow{4}{4.0cm}{Are the scientific claims factually correct?}
& 1--3 & Major factual errors that undermine the answer. \\
& & 4--5 & Mostly correct with noticeable inaccuracies. \\
& & 6--7 & Correct with minor imprecisions. \\
& & 8--10 & Fully correct; all claims are verifiable and accurate. \\
\midrule

% relevance
\multirow{4}{2.6cm}{Relevance}
& \multirow{4}{4.0cm}{Does the answer address what was asked?}
& 1--3 & Off-topic; does not address the question. \\
& & 4--5 & Partially addresses the question. \\
& & 6--7 & Addresses the question with minor tangential content. \\
& & 8--10 & Precisely targeted to the question. \\
\midrule

% completeness
\multirow{4}{2.6cm}{Completeness}
& \multirow{4}{4.0cm}{Are all aspects of the question covered?}
& 1--3 & Major gaps; key aspects are missing. \\
& & 4--5 & Covers some aspects but omits important ones. \\
& & 6--7 & Covers most aspects with minor omissions. \\
& & 8--10 & Comprehensive; all aspects addressed. \\
\midrule

% directness
\multirow{4}{2.6cm}{Directness}
& \multirow{4}{4.0cm}{Is the key point stated upfront or buried?}
& 1--3 & Key point buried after 200$+$ words of preamble. \\
& & 4--5 & Takes too long to reach the core answer. \\
& & 6--7 & Reasonably direct with moderate lead-in. \\
& & 8--10 & Key answer stated clearly within the first few sentences. \\
\midrule

% quantitative precision
\multirow{4}{2.6cm}{Quantitative Precision}
& \multirow{4}{4.0cm}{Are specific numbers included where relevant?}
& 1--3 & No quantitative details despite relevance. \\
& & 4--5 & Vague references to numbers (e.g.\ ``some studies show\ldots''). \\
& & 6--7 & Some specific values provided. \\
& & 8--10 & Rich in quantitative detail (rates, ranges, dates, areas). \\
\midrule

% response structure
\multirow{4}{2.6cm}{Response Structure}
& \multirow{4}{4.0cm}{Is the format appropriate for the content?}
& 1--3 & Wall of text; no structure. \\
& & 4--5 & Basic paragraphs with minimal organisation. \\
& & 6--7 & Decent organisation with headings or bullets where helpful. \\
& & 8--10 & Optimal structure for the question type. \\
\midrule

% scope awareness
\multirow{4}{2.6cm}{Scope Awareness}
& \multirow{4}{4.0cm}{Are limitations and boundaries acknowledged?}
& 1--3 & No qualification; presents universal claims without caveats. \\
& & 4--5 & Minimal acknowledgement of scope limits. \\
& & 6--7 & Some caveats mentioned. \\
& & 8--10 & Explicitly flags data gaps, regional limits, or temporal bounds. \\
\midrule

% helpfulness
\multirow{4}{2.6cm}{Helpfulness}
& \multirow{4}{4.0cm}{Is the answer actionable for researchers?}
& 1--3 & Not useful; vague or misleading. \\
& & 4--5 & Provides a basic overview only. \\
& & 6--7 & Useful with some actionable information. \\
& & 8--10 & Highly actionable; clearly advances the reader's understanding. \\
\bottomrule
\end{longtable}

\chapter{Initial Evaluation Criteria}\label{app:initial-criteria}

Table~\ref{tab:initial-criteria-full} lists the five criteria used in
the initial GPT-4.1-mini evaluation framework described in
\S\ref{sec:initial-framework}. These criteria were modelled after the
evaluation framework of~\citet{elbaff2025earthobservationrag} to ensure
comparability with prior results. Each criterion was scored on a 1--5
integer scale.

\begin{table}[htbp]
\centering
\caption{Criteria and definitions used in the initial GPT-4.1-mini evaluation framework.}
\label{tab:initial-criteria-full}
\begin{tabularx}{\textwidth}{lX}
\toprule
\textbf{Criterion} & \textbf{Definition} \\
\midrule
Factuality & Accuracy of scientific claims made in the answer. \\
Relevance & Alignment of the answer with the question's intent. \\
Groundedness & Degree to which claims are supported by the provided context and sources. \\
Helpfulness & Practical value of the answer for a domain researcher. \\
Depth & Level of technical detail and completeness of the response. \\
\bottomrule
\end{tabularx}
\end{table}

As discussed in Chapter~\ref{chap:evaluation}, this configuration
produced score ceiling saturation (93\% of DLR two-step scores and 99\%
of ontology-enhanced scores at the maximum of~5) and lacked criteria
targeting the four failure patterns identified in
Chapter~\ref{chap:preliminaries}. These limitations motivated the
transition to the eight-criterion Claude Sonnet~4.6 framework
(Appendix~\ref{app:eval-criteria}).

\chapter{Agent Prompts}\label{app:prompts}

This appendix reproduces the full system prompts used by the three agents.
Template placeholders written in braces (e.g.\ \texttt{\{question\}}) are filled at
runtime. The prompts are referenced from Chapter~\ref{chap:approach}; the design
choices behind each prompt are discussed in the corresponding Prompt Design
subsections there.

\section{Ontology Agent Extraction Prompt}\label{app:prompt-ontology}

Used by the Ontology Agent to extract 4--8 attribute-value pairs from the user question.
Referenced in Section~\ref{subsec:ontology-prompt-design}.

\begin{lstlisting}[style=promptstyle,
caption={Ontology Agent extraction prompt (full).},
label={lst:ontology-prompt}]
TASK: Extract semantic attributes from Earth Observation questions

Extract 4-8 attribute-value pairs capturing key concepts essential for answering the question.
Focus on semantic completeness - cover all major scientific dimensions mentioned.

MANDATORY CATEGORIES (cover >=4 of 5):
1. BIOPHYSICAL_VARIABLE: What's measured? (glacier_retreat, sea_ice_extent, albedo, carbon_flux)
2. SENSOR/METHOD: How measured? (Sentinel-2, MODIS, SAR, time_series_analysis)
3. SPATIAL_QUALIFIER: Where? (global, polar_regions, 1km_resolution)
4. TEMPORAL_QUALIFIER: When? (decadal_trends, 1990-present, seasonal_cycle)
5. APPLICATION_DOMAIN: Why? (climate_assessment, ecosystem_monitoring)

CENTRALITY SCALE (6-point):
- 1.0: ABSOLUTELY CRITICAL (answer impossible without this)
- 0.9: CORE MECHANISM (primary driver of answer)
- 0.7: IMPORTANT ENABLER (significantly constrains/enables)
- 0.5: CONTEXTUAL SUPPORT (enriches but not essential)
- 0.3: BACKGROUND FRAMEWORK (contextual framing)
- 0.1: PERIPHERAL (minimal value)

QUESTION
{question}

SEMANTIC COMPLETENESS GUIDANCE (CRITICAL)

Focus on COMPLETE semantic coverage, NOT arbitrary pair counts.
Extract ALL distinct concepts needed to fully answer the question.

COMPLEXITY TIERS (pair counts are GUIDELINES, not targets):

Simple questions (typically 4-5 pairs):
Example: "What is albedo?"
- Core concept (1.0): albedo
- Measurement approach (0.9): satellite_observation
- Temporal context (0.7): seasonal_variation
- Application (0.5): climate_modeling

Standard questions (typically 5-6 pairs):
Example: "How do glaciers influence climate?"
- Core mechanisms (1.0/0.9): glacier_retreat, albedo_feedback
- Observational methods (0.7): satellite_altimetry
- Temporal dynamics (0.7): long_term_trends
- Spatial context (0.5): polar_regions
- Applications (0.5): climate_change_assessment

Complex questions (typically 6-8 pairs):
Example: "How do variations in glacier retreat, sea ice extent, and albedo influence carbon flux?"
- Multiple core variables (1.0/0.9): glacier_retreat, sea_ice_extent, albedo, carbon_flux
- Interaction mechanisms (0.9): ice_albedo_feedback, carbon_cycle_dynamics
- Measurement approaches (0.7): multi_sensor_integration
- Temporal dimensions (0.7): decadal_trends, seasonal_cycles
- Spatial dimensions (0.5): Arctic_vs_Antarctic, regional_variability
- Applications (0.5): climate_projections, ecosystem_impacts

SEMANTIC VALIDATION CHECKLIST:
[x] Does every major scientific concept in the question have a corresponding attribute?
[x] Are all implied measurement/observation methods captured?
[x] Are geographic and temporal constraints properly represented?
[x] Is the scientific purpose/application context included?
[x] Can a domain expert answer the question using ONLY these pairs?

VALUE SPECIFICITY: Avoid vague or generic values.
Good: glacier_velocity, surface_albedo_change, CH4_flux, NDVI_anomaly, ENSO_index
Bad: change, data, context, conditions

INSTRUCTIONS:
1. Read the QUESTION, identifying all distinct concepts and mechanisms.
2. Assess question complexity: simple / standard / complex.
3. Extract pairs accordingly. Do NOT pad or truncate.
4. Ensure you cover >=4 of the 5 mandatory categories.
5. Assign centrality based on question structure (honest distribution).
6. Each description must be 1-2 sentences, specific to the question context.

RETURN ONLY THIS JSON STRUCTURE (valid JSON, no markdown):

"pairs": [

"attribute": "attribute_name",
"value": "specific_value",
"description": "what this means in the context of the question",
"value_type": "categorical",
"constraints": [],
"centrality": 1.0

\end{lstlisting}

\section{Refinement Agent Prompt (Ontology-Enhanced Variant)}\label{app:prompt-refinement}

Used by the Refinement Agent when an ontology is available. The \texttt{\{ontology\}} placeholder
is filled with the sorted attribute-value pairs from the Ontology Agent, ordered by centrality
in descending order. Referenced in Section~\ref{subsec:refinement-prompt-design}.

\begin{lstlisting}[style=promptstyle,
caption={Refinement Agent prompt -- ontology-enhanced variant (full).},
label={lst:refinement-prompt}]
You are a research context refinement assistant.

Goal: From the AQL results and ontology, build a compact, well-structured TEXT context that a
downstream answer model can use to answer ONE user question. The context must stay within the
ontology and AQL evidence, be easy to scan, and avoid placeholders.

================================================================
INPUT
- USER QUESTION:
{query}

- ONTOLOGY (target attributes / concepts):
{ontology}

AQL RESULTS are numbered; scan them from 1 to N and use only the 3-8 most ontology-relevant
items for the context.
When selecting the 3-8 most relevant AQL RESULTS:
- Prefer items whose title or first 3 sentences explicitly mention the key ONTOLOGY attributes
in {ontology} (e.g., main biophysical variables or application domains).
- If more than 8 items match, choose the 3-8 most recent by publication year.
- If fewer than 3 items match, use all matching items and state the limited coverage in
[CAVEATS & GAPS].

- AQL RESULTS (primary evidence: data + document metadata):
{aql_results}
================================================================

GENERAL RULES
- Use ONLY information from the ONTOLOGY and AQL_RESULTS. Do NOT use any training knowledge.
- Hallucination guardrail:
- Do NOT invent numbers, regions, years, methods, or citations.
- If something is missing, write "Not available in AQL results" instead of a placeholder.
- Do NOT create new references or modify any reference metadata.
- Be concise but informative. Include key methodological details (sensor, algorithm, resolution)
and at least one short explanatory sentence per topic when available in AQL_RESULTS.
- Always specify scope (region + time) when giving quantitative results.
- Clearly separate regional vs global findings.
- For each critical attribute, if the AQL_RESULTS provide methods or drivers (e.g., satellite
sensor, retrieval algorithm, forcing factors), append a brief clause describing them.

================================================================
OUTPUT FORMAT (plain text, no JSON, no code fences)
Include EACH of the following sections EXACTLY ONCE and in this ORDER.

[ONTOLOGY SUMMARY]
Critical attributes found:
- attribute_name: value (with units, region, time, method if relevant)
- ...

Supporting attributes found:
- attribute_name: value
- ...

[TOPICS AND INFORMATION]
# Topic name 1
Findings: 1-2 key findings for this topic (quantitative and/or qualitative).
Evidence: Short citation from AQL (e.g., "Author et al., 2020").
Mechanisms / methods: Brief description of the measurement approach or main drivers.
Scope: Geographic: {region or "global"} | Temporal: {years}.
Confidence: {0-1} with 1 short reason.
Limitations: 1-2 specific caveats.
Use: Short instruction for the answer model.

[VALIDATED REFERENCES]
{documents}

For each reference included in the final context:
1. Author et al. (Year). *Title*. URI.
Scope: Geographic: Region | Temporal: Period
Key contribution: Main finding.
Limitations: Key limitations.
Use: How to use this reference.

DO NOT INVENT OR MODIFY ANY REFERENCE INFORMATION.
ONLY INCLUDE REFERENCES DIRECTLY USEFUL TO ANSWER THE QUESTION.

RE-SCAN FOCUS: Prioritize matches to QUERY {query} and ONTOLOGY {ontology}

[CAVEATS & GAPS]
Ontology coverage:
- Addressed: {list key ontology attributes well covered}.
- Missing or weak: {attributes not covered or poorly supported}.

Data coverage:
- Geographic: Covered: {regions} | Gaps: {regions}.
- Temporal: Covered: {periods} | Gaps: {periods}.

Key limitations:
- 1-3 bullets with the most important limitations for answering the question.

Generation guidance:
- Emphasise: 2-3 concepts or metrics best supported and most relevant to the question.
- Caveat: 1-2 concepts that must always be qualified when mentioned.
- Frame answer as: 1 sentence high-level framing suggestion for the final answer model.

================================================================
ADDITIONAL CONSTRAINTS
- Replace all template braces {...} with actual values or short phrases.
- Prefer specific terms for geographic and temporal references.
- Organise topics by concept (matching the ontology attributes), not by document.
- Do not output any evaluation of the user, model, or system; only the context.
\end{lstlisting}

\section{Refinement Agent Prompt (No-Ontology Variant)}\label{app:prompt-refinement-no-ontology}

Used by the Refinement Agent in the no-ontology ablation variant (Section~\ref{subsec:refinement-prompt-design}).
The \texttt{ONTOLOGY} input block is removed and the document selection criterion changes from
``prefer items that mention key ONTOLOGY attributes'' to ``prefer items directly relevant to
the user question.'' The \texttt{RE-SCAN FOCUS} line no longer includes the ontology anchor.
The output format is otherwise identical, enabling a clean ablation.

\begin{lstlisting}[style=promptstyle,
caption={Refinement Agent prompt -- no-ontology ablation variant (full).},
label={lst:refinement-prompt-no-ontology}]
You are a research context refinement assistant.

Goal: From the AQL results, build a compact, well-structured TEXT context that a downstream
answer model can use to answer ONE user question. The context must stay within the AQL evidence,
be easy to scan, and avoid placeholders.

NOTE: This experiment does not use an ontology. Select and organise AQL results based solely
on their relevance to the user question.

================================================================
INPUT
- USER QUESTION:
{query}

AQL RESULTS are numbered; scan them from 1 to N and use only the 3-8 most relevant items.
When selecting the 3-8 most relevant AQL RESULTS:
- Prefer items whose title or first 3 sentences are directly relevant to the user question.
- If more than 8 items match, choose the 3-8 most recent by publication year.
- If fewer than 3 items match, use all matching items and state the limited coverage in
[CAVEATS & GAPS].

- AQL RESULTS (primary evidence: data + document metadata):
{aql_results}
================================================================

GENERAL RULES
- Use ONLY information from the AQL_RESULTS. Do NOT use any training knowledge.
- Hallucination guardrail:
- Do NOT invent numbers, regions, years, methods, or citations.
- If something is missing, write "Not available in AQL results" instead of a placeholder.
- Do NOT create new references or modify any reference metadata.
- Be concise but informative. Use bullets, but include key methodological details
(sensor, algorithm, resolution) and at least one short explanatory sentence per topic.
- Always specify scope (region + time) when giving quantitative results.
- Clearly separate regional vs global findings.

================================================================
OUTPUT FORMAT (plain text, no JSON, no code fences)
Include EACH of the following sections EXACTLY ONCE and in this ORDER.

[TOPICS AND INFORMATION]
# Topic name 1
Findings: 1-2 key findings for this topic (quantitative and/or qualitative).
Evidence: Short citation from AQL (e.g., "Author et al., 2020").
Mechanisms / methods: Brief description of the measurement approach or main drivers.
Scope: Geographic: {region or "global"} | Temporal: {years}.
Confidence: {0-1} with 1 short reason.
Limitations: 1-2 specific caveats.
Use: Short instruction for the answer model.

[VALIDATED REFERENCES]
{documents}

For each reference included in the final context:
1. Author et al. (Year). *Title*. URI.
Scope: Geographic: Region | Temporal: Period
Key contribution: Main finding.
Limitations: Key limitations.
Use: How to use this reference.

DO NOT INVENT OR MODIFY ANY REFERENCE INFORMATION.
ONLY INCLUDE REFERENCES DIRECTLY USEFUL TO ANSWER THE QUESTION.

RE-SCAN FOCUS: Prioritize matches to QUERY {query}

[CAVEATS & GAPS]
Data coverage:
- Geographic: Covered: {regions} | Gaps: {regions}.
- Temporal: Covered: {periods} | Gaps: {periods}.

Key limitations:
- 1-3 bullets with the most important limitations for answering the question.

Generation guidance:
- Emphasise: 2-3 concepts or metrics best supported and most relevant to the question.
- Caveat: 1-2 concepts that must always be qualified when mentioned.
- Frame answer as: 1 sentence high-level framing suggestion for the final answer model.

================================================================
ADDITIONAL CONSTRAINTS
- Replace all template braces {...} with actual values or short phrases.
- Prefer specific terms for geographic and temporal references.
- Organise topics by concept, not by document.
- Do not output any evaluation of the user, model, or system; only the context.
\end{lstlisting}

\section{Generation Agent Refinement Prompt}\label{app:prompt-generation}

Used in Step~2 of the Generation Agent. The \texttt{\{draft\_answer\}} placeholder is filled with
the zero-shot draft from Step~1; \texttt{\{context\}} is filled with the \texttt{enriched\_context}
output of the Refinement Agent. Referenced in Section~\ref{subsec:generation-prompt-design}.

\begin{lstlisting}[style=promptstyle,
caption={Generation Agent refinement prompt (full).},
label={lst:generation-prompt}]
You are an expert assistant specialized in Earth Observation (EO) science. Your task is to
answer scientific questions in clear, research-appropriate language.

Most important rules (follow these even if other instructions conflict):
- Base all key claims primarily on the provided CONTEXT when possible.
- If the draft answer and CONTEXT conflict, prefer the CONTEXT for numbers, regions,
and time windows.
- Do not invent or modify quantitative values (years, rates, areas) that contradict
the CONTEXT.
- Keep the answer structured around the question.
- When the CONTEXT provides sufficient evidence, give a deep answer: explain mechanisms,
methods, and uncertainties step-by-step, instead of only listing findings. Prefer 2-3
short reasoning steps over a single high-level sentence.
- Match the complexity of the question to the depth and detail of the answer.
- Explicitly mention major uncertainties or data gaps when they appear in the CONTEXT.
- Use at least one numeric citation marker (e.g., [1]) when the CONTEXT provides
VALIDATED REFERENCES, and include a matching References section.

You answered a question with the following draft answer, using your knowledge.
Now refine this answer using the provided context.

### YOUR DRAFT ANSWER
- Question: {question}
- Draft Answer: {draft_answer}

### CONTEXT
The context contains publications and their related topics extracted from an Earth
Observation knowledge graph:

{context}

RE-SCAN FOCUS: Remember the QUERY {question}

### YOUR ROLE
Treat the draft answer as a strong starting point and the context as high-priority
additional evidence:
- Keep what is correct and useful in the draft.
- Use the context to correct, sharpen, or extend the draft.
- You MAY use your training knowledge to add relevant details, as long as they do not
conflict with the CONTEXT and are clearly standard background.

### PRIORITY RULES
1. Decide whether the question is GENERAL (global/conceptual) or SPECIFIC (focused on
a particular region, system, dataset, or numerical estimate).

2. If the draft and context AGREE on a fact:
- Keep it, and you may enrich it with extra detail.
- Prioritise the 3-5 most important facts; you do NOT need to use every detail.

3. If the draft and context CONFLICT on a specific fact:
- Align the draft with the CONTEXT by correcting conflicting numbers, regions,
or time windows.

4. If something important is in the context but missing from the draft:
- Add it only if it clearly helps answer the question.

5. If something is in the draft but NOT in the context:
- Keep it if it is standard, well-established background AND does not contradict
the context.
- Otherwise, remove it or mark it clearly as a tentative interpretation.

6. When multiple CONTEXT items report related findings, synthesise them explicitly
(e.g., "Several studies show...") rather than describing each one separately.
Support each synthesis with one or more numeric citations when backed by references.

### STRUCTURE
- Begin with a short direct answer (2-4 sentences) that clearly addresses the question.
- Then add 2-4 short sections with headings that match the question type:
"Key concepts" / "Definition" for "what is..." questions
"Main mechanisms" / "Impacts" for "how/why..." questions
"Comparison of X and Y" for comparison questions
"Strengths" / "Limitations" for evaluation questions
- For comparison questions, include ONE small comparison table if it helps clarity.
- When the question asks about mechanisms or methodology, add a "Detailed mechanisms
and methods" section (2-4 sentences).
- If the context includes important caveats or confidence information, add a short
"Uncertainties and limitations" paragraph (2-4 sentences).
- When there is clear real-world or modelling relevance, end with a brief "Implications"
paragraph (2-4 sentences).

### USE OF NUMBERS AND SCOPE
- Use quantitative details from the CONTEXT and your training knowledge, as long as
they are consistent.
- Whenever you give a quantitative value from the CONTEXT, mention:
* the geographic scope (e.g., "Hindu Kush Himalaya, 27-37N, 66-92E")
* the temporal scope (e.g., "2000-2024")
- If the question is global or very general:
* Keep the main narrative GENERAL.
* Use regional details from the CONTEXT only as BRIEF examples, clearly labelled.
* Use at most one short paragraph per example region.
- If the question is explicitly regional:
* Focus the answer on that region.
* Use global behaviour only as supporting context.

### UNCERTAINTY AND CAVEATS
- Reflect major data gaps (e.g., regions or periods not covered).
- Use calibrated language ("high confidence", "medium confidence", "limited evidence").
- Highlight only the 2-3 most important limitations.

### STYLE AND FORMAT
- Write in clear, neutral, research-appropriate language.
- Be informative but not verbose; avoid repeating the same numbers or phrases.
- Do NOT copy internal context headers or markers such as "[ONTOLOGY SUMMARY]",
"[TOPICS AND INFORMATION]", "[CAVEATS & GAPS]", or any bracketed labels from context.
- The ONLY bracketed items you may output are numeric citation markers like [1], [2].
- Do NOT mention AQL, knowledge graphs, prompts, or any internal pipeline details.
- Provide ONLY the final refined answer, with no meta-commentary.

### CITATIONS
- The CONTEXT includes a [VALIDATED REFERENCES] section with formatted references.
- Insert numeric markers [1], [2] in the text for key claims.
- Add a "References" section at the end listing used references in the same format
as provided.
- Do NOT fabricate or hallucinate references.
- Use the exact format provided in the [VALIDATED REFERENCES] section.

Now write the final refined answer.
\end{lstlisting}

\chapter{JSON Schemas}\label{app:schemas}

This appendix provides the JSON schemas for the structured data objects
passed between pipeline agents. The schemas are referenced in
Chapter~\ref{chap:approach}.

\section{Dynamic Ontology Schema}\label{app:schema-ontology}

\begin{lstlisting}[style=promptstyle,
caption={JSON schema for the \texttt{DynamicOntology} object produced
by the Ontology Agent.},
label={lst:schema-ontology}]
{
  "pairs": [
    {
      "attribute": "string (category name)",
      "value": "string (specific concept)",
      "description": "string (1-2 sentences contextualising the pair)",
      "value_type": "categorical | numerical | temporal | spatial",
      "constraints": ["string (units, ranges, or properties)"],
      "centrality": "float (one of: 0.1, 0.3, 0.5, 0.7, 0.9, 1.0)"
    }
  ]
}
\end{lstlisting}

\section{Refined Context Schema}\label{app:schema-context}

\begin{lstlisting}[style=promptstyle,
caption={Structure of the \texttt{RefinedContext} object produced by
the Refinement Agent and consumed by the Generation Agent.},
label={lst:schema-context}]
RefinedContext:
  original_context: string    # raw AQL results (preserved for audit)
  question: string            # original user question
  assessments: list           # minimal per-document metadata
  enriched_context: string    # structured plain-text context
    Sections:
      [ONTOLOGY SUMMARY]      # critical and supporting attributes
      [TOPICS AND INFORMATION]# per-topic findings, evidence, scope,
                              # confidence, limitations, usage guidance
      [VALIDATED REFERENCES]  # bibliographic metadata from OpenAlex
      [CAVEATS & GAPS]        # ontology and data coverage analysis
\end{lstlisting}

\chapter{Agent Code Excerpts}\label{app:code}

The following excerpts highlight the core logic of each agent.
Full source code is available in the project repository.

\section{Ontology Agent}\label{app:code-ontology}

\begin{lstlisting}[style=pythonstyle,
caption={Ontology Agent: \texttt{process} and
\texttt{\_extract\_attribute\_value\_pairs} methods.},
label={lst:code-ontology}]
class OntologyConstructionAgent(BaseAgent):
    # Constructs a dynamic ontology from a user question via LLM calls.

    def process(self, question: str,
                include_relationships: bool = True) -> DynamicOntology:
        av_pairs = self._extract_attribute_value_pairs(question)
        relationships = []
        if include_relationships:
            relationships = self._define_logical_relationships(av_pairs, question)
        return DynamicOntology(
            attribute_value_pairs=av_pairs,
            logical_relationships=relationships,
            source_query=question,
            should_use_ontology=True,
        )

    def _extract_attribute_value_pairs(self, question: str) -> List[AttributeValuePair]:
        template = self.load_prompt("extract_attributes_slim.txt")
        prompt = template.replace("{question}", question)
        response = self.llm.invoke(prompt, force_json=True)
        if isinstance(response, dict) and "pairs" in response:
            return [
                AttributeValuePair(
                    attribute=p["attribute"],
                    value=p["value"],
                    description=p.get("description", ""),
                    centrality=p.get("centrality", 0.5),
                )
                for p in response["pairs"]
            ]
        return []
\end{lstlisting}

\section{Refinement Agent}\label{app:code-refinement}

\begin{lstlisting}[style=pythonstyle,
caption={Refinement Agent: \texttt{process\_context} and
\texttt{\_build\_refined\_prompt} methods.},
label={lst:code-refinement}]
class RefinementAgent1PassRefined(BaseRefinementAgent):
    # Single-pass refinement: sends AQL results + ontology to the LLM
    # and receives a ready-to-use structured text context.

    def process_context(self, question, structured_context,
                        ontology=None, aql_results_str=None):
        if not aql_results_str:
            return self._empty_context(question, ontology)
        prompt, documents = self._build_refined_prompt(
            question, ontology, include_ontology=True,
            aql_results_str=aql_results_str
        )
        refined_text = self.llm.invoke(prompt).strip()
        return RefinedContext(
            original_context=structured_context,
            question=question,
            assessments=self._minimal_assessments(documents),
            enriched_context=refined_text,
        )

    def _build_refined_prompt(self, question, ontology,
                              include_ontology, aql_results_str):
        template = self.load_prompt("refinement_1pass_refined_exp4.txt")
        ontology_text = (
            self._format_ontology(ontology)
            if include_ontology and ontology else "No ontology provided"
        )
        aql_text, documents = self._parse_aql_for_prompt(aql_results_str)
        documents_text = self._format_documents_for_references(documents)
        prompt = (
            template
            .replace("{query}", question)
            .replace("{ontology}", ontology_text)
            .replace("{aql_results}", aql_text)
            .replace("{documents}", documents_text)
        )
        return prompt, documents
\end{lstlisting}

\section{Generation Agent}\label{app:code-generation}

\begin{lstlisting}[style=pythonstyle,
caption={Generation Agent: \texttt{generate} and
\texttt{\_build\_refinement\_prompt} methods.},
label={lst:code-generation}]
class GenerationAgent(BaseAgent):
    # Two-step generation: zero-shot draft then context-grounded refinement.

    def generate(self, question, text_context, ontology=None):
        # Step 1: zero-shot draft (question only)
        zero_prompt = self.zero_shot_template.replace("{question}", question)
        draft = self._call_llm(zero_prompt)
        # Step 2: refine using retrieved context
        refine_prompt = self._build_refinement_prompt(question, draft, text_context, ontology)
        final_answer = self._call_llm(refine_prompt)
        return Answer(
            answer=final_answer,
            formatted_references=self._extract_formatted_references(final_answer),
        )

    def _build_refinement_prompt(self, question, draft, context, ontology):
        prompt = self.refinement_template
        prompt = prompt.replace("{question}", question)
        prompt = prompt.replace("{draft_answer}", draft)
        prompt = prompt.replace(
            "{context}", context.strip() if context else "No context available."
        )
        if ontology and ontology.attribute_value_pairs:
            ont_lines = [
                f"- {av.attribute}: {av.value} ({av.description})"
                for av in ontology.attribute_value_pairs
            ]
            prompt = prompt.replace("{ontology}", "\n".join(ont_lines))
        else:
            prompt = prompt.replace("{ontology}", "No ontology constraints.")
        return prompt
\end{lstlisting}
